\SourceLeaf{362}
adeoque valores integralium
\[
\begin{gathered}\int\frac{\mathrm{d}T}{2\pi\surd(mm\cos T^2+nn\sin T^2)},\\[7pt]
\int\frac{\mathrm{d}T'}{2\pi\surd(m'm'\cos T'^2+n'n'\sin T'^2)}\end{gathered}
\]
si utriusque variabilis a valore $0$ usque ad valorem $360^{0}$ extenditur, inter se aequales. Et quum perinde ulterius continuare liceat, patet, his valoribus etiam
aequalem esse valorem integralis
\[
\int\frac{\mathrm{d}\Theta}{2\pi\surd(\mu\mu\cos\Theta^2+\mu\mu\sin\Theta^2)}
\]
a $\Theta=0$ usque ad $\Theta=360^{0}$, qui manifesto fit $=\frac{1}{\mu}$. Q.~E.~D.

\AttractionArticle{17}
Ex aequatione, relationem inter $T$ et $T'$ exhibente,
\[
(m-n)\sin T.\sin T'^2=2m\sin T'-(m+n)\sin T
\]
facile deducitur
\[
\begin{aligned}
\surd(mm\cos T^2+nn\sin T^2)&=m-(m-n)\sin T.\sin T'\\[6pt]
\surd(m'm'\cos T'^2+n'n'\sin T'^2)&=m\mathop{\mathrm{cotang}}T.\mathop{\mathrm{tang}}T'
\end{aligned}
\]
atque hinc, adiumento eiusdem aequationis,
\[
\begin{aligned}
&\sin T.\sin T'.\surd(mm\cos T^2+nn\sin T^2)+m'(\cos T^2-\sin T^2)\\[6pt]
&=\cos T.\cos T'.\surd(m'm'\cos T'^2+n'n'\sin T'^2)-\tfrac12(m-n)\sin T'^2
\end{aligned}
\]
\par Multiplicata hac aequatione per
\[
\frac{\mathrm{d}T}{\surd(mm\cos T^2+nn\sin T^2)}
=\frac{\mathrm{d}T'}{\surd(m'm'\cos T'^2+n'n'\sin T'^2)}
\]
prodit
\[
\begin{aligned}&\frac{m'(\cos T^2-\sin T^2)\mathrm{d}T}{\surd(mm\cos T^2+nn\sin T^2)}\\[7pt]&=-\frac{\tfrac12(m-n)\sin T'^2\mathrm{d}T'}{\surd(m'm'\cos T'^2+n'n'\sin T'^2)}+\mathrm{d}.\sin T'\cos T\end{aligned}
\]
\par Multiplicando hanc aequationem per $\frac{m-n}{\pi}$, substituendo
$m'(m-n)=\tfrac12(mm-nn)$,
$(m-n)^2=4(m'm'-n'n')$,
$\sin T'^2=\tfrac12-\tfrac12(\cos T'^2-\sin T'^2)$, et integrando, a valoribus $T$ et $T'=0$ usque ad $360^{0}$, habemus:
\[
\begin{aligned}
&(mm-nn)\int\frac{(\cos T^2-\sin T^2).\mathrm{d}T}{2\pi\surd(mm\cos T^2+nn\sin T^2)}\\[6pt]
&=-\frac{2(m'm'-n'n')}{\mu}\\[7pt]&\quad+2(m'm'-n'n')\int\frac{(\cos T'^2-\sin T'^2)\mathrm{d}T'}{2\pi\surd(m'm'\cos T'^2+n'n'\sin T'^2)}
\end{aligned}
\]
\SourceLeaf{363}
Et quum integrale definitum ad dextram perinde transformare liceat, manifesto
integrale
\[
\int\frac{(\cos T^2-\sin T^2)\mathrm{d}T}{2\pi\surd(mm\cos T^2+nn\sin T^2)}
\]
exprimetur per seriem infinitam citissime convergentem
\[
-\frac{\begin{gathered}2(m'm'-n'n')+4(m''m''-n''n'')\\+8(m'''m'''-n'''n''')+\mathrm{etc.}\end{gathered}}{(mm-nn)\mu}
=-\frac{\nu}{\mu}
\]
\par Calculus numericus commodissime per logarithmos perficitur, si statuimus
\[
\begin{gathered}\tfrac14\surd(mm-nn)=\lambda,\quad\tfrac14\surd(m'm'-n'n')=\lambda',\\[6pt]\tfrac14\surd(m''m''-n''n'')=\lambda''\ \mathrm{etc.}\end{gathered}
\]
unde erit
\[
\lambda'=\frac{\lambda\lambda}{m'},\quad
\lambda''=\frac{\lambda'\lambda'}{m''},\quad
\lambda'''=\frac{\lambda''\lambda''}{m'''}\ \mathrm{etc.}\quad\text{atque}
\]
\[
\nu=\frac{2\lambda'\lambda'+4\lambda''\lambda''+8\lambda'''\lambda'''+\mathrm{etc.}}{\lambda\lambda}
\]

\AttractionArticle{18}
Per methodum hic explicatam etiam integralia \textit{indefinita} (a valore variabilis
$=0$ inchoantia) maxima concinnitate assignare licet. Scilicet, si $T''$ perinde
per $m',\ n',\ T'$ determinari supponitur, uti $T'$ per $m.\ n,\ T$, ac perinde rursus
$T'''$ per $m'',\ n'',\ T''$ et sic porro, etiam pro quovis valore determinato ipsius $T$,
valores terminorum serie $T,\ T',\ T'',\ T'''$ etc. ad limitem $\Theta$ citissime convergent,
eritque
\[
\int\frac{\mathrm{d}T}{\surd(mm\cos T^2+nn\sin T^2)}=\frac{\Theta}{\mu}
\]
\[
\begin{aligned}&\int\frac{(\cos T^2-\sin T^2)\mathrm{d}T}{\surd(mm\cos T^2+nn\sin T^2)}\\[7pt]&=-\frac{\nu\Theta}{\mu}\\[7pt]&\quad+\frac{\begin{gathered}\lambda'\cos T\sin T'+2\lambda''\cos T'\sin T''\\+4\lambda'''\cos T'\sin T'''+\mathrm{etc.}\end{gathered}}{\lambda\lambda}\end{aligned}
\]
\par Sed haec obiter hic addigitavisse sufficiat, quum ad institutum nostrum non
sint necessaria.

\AttractionArticle{19}
Quodsi iam statuimus $m=\surd(G+G')$, $n=\surd(G+G'')$, valores quantitatum $P,\ Q$ facile ad transscendentes $\mu,\ \nu$ reducentur. Quum enim $P,\ Q$ sint
valores integralium
\SourceLeaf{364}
\[
\begin{gathered}\int\frac{\cos T^2\mathrm{d}T}{2\pi(mm\cos T^2+nn\sin T^2)^{\frac32}},\\[7pt]\int\frac{\sin T^2\mathrm{d}T}{2\pi(mm\cos T^2+nn\sin T^2)^{\frac32}}\end{gathered}
\]
a $T=0$ usque ad $T=360^{0}$ extensorum, primo statim obvium est, haberi
\[
mmP+nnQ=\frac{1}{\mu}\quad .\, .\, .\, .\, .\, .\, .\, .\, .\, .\, .\, .\quad[24]
\]
\par Porro fit
\[
\begin{aligned}
&\frac{(\cos T^2-\sin T^2)\mathrm{d}T}{2\pi\surd(mm\cos T^2+nn\sin T^2)}\\[7pt]&\quad+\frac{(mm\cos T^2-nn\sin T^2)\mathrm{d}T}{2\pi(mm\cos T^2+nn\sin T^2)^{\frac32}}\\[6pt]
&\qquad=\frac{(mm\cos T^4-nn\sin T^4)\mathrm{d}T}{\pi(mm\cos T^2+nn\sin T^2)^{\frac32}}\\[6pt]
&\qquad=\mathrm{d}.\frac{\cos T\sin T}{\pi\surd(mm\cos T^2+nn\sin T^2)}
\end{aligned}
\]
\par Integrando hanc aequationem a $T=0$ usque ad $T=360^{0}$, prodit
\[
-\frac{\nu}{\mu}+mmP-nnQ=0\quad .\, .\, .\, .\, .\, .\, .\, .\, .\quad[25]
\]
\par E combinatione aequationum 24, 25 denique colligimus
\[
P=\frac{1+\nu}{2mm\mu},\qquad Q=\frac{1-\nu}{2nn\mu}
\]
\begin{center}\rule{0.22\textwidth}{0.4pt}\end{center}
\clearpage\fancyhead[L]{\small ANZEIGE. DETERMINATIO ATTRACTIONIS}
\pdfbookmark[0]{ANZEIGE. DETERMINATIO ATTRACTIONIS}{work-native366}
\SourceLeaf{366}
\selectlanguage{german}
\begin{center}
\textbf{ANZEIGE.}
\end{center}
\begin{center}
\rule{0.62\linewidth}{0.4pt}\\[6pt]
Göttingische gelehrte Anzeigen.\quad 1818 Februar 9.\\[6pt]
\rule{0.62\linewidth}{0.4pt}
\end{center}

Am $17^{\mathrm{ten}}$ Januar übergab Hr. Hofr. \textsc{Gauss} der Königl. Societät eine Vorlesung:

\begin{center}
\textit{Determinatio attractionis, quam in punctum quodlibet positionis datae exerceret
planeta, cuius massa per totam eius orbitam, ratione temporis, quo singulae partes
describuntur, uniformiter esset dispertita.}
\end{center}

Vermöge eines, vielleicht bis jetzt noch von niemand ausdrücklich ausgesprochenen, aber aus den Gründen der physischen Astronomie leicht zu beweisenden Lehrsatzes, sind die Säcularveränderungen einer Planetenbahn durch die
Störung eines andern Planeten dieselben, der störende Planet mag seine elliptische Bahn nach Keplers Gesetzen wirklich beschreiben, oder seine Masse mag
auf den Umfang der Ellipse in dem Masse vertheilt angenommen werden, dass
auf Stücke der Ellipse, die sonst in gleich grossen Zeiten beschrieben werden.
gleich grosse Antheile an der ganzen Masse kommen: vorausgesetzt, dass die Umlaufszeiten des gestörten und des störenden Planeten nicht in rationalem Verhältnisse zu einander stehen. Die Aufgabe, welche den Gegenstand dieser Abhandlung ausmacht, nemlich die nicht genäherte, sondern genaue, nicht von mässiger Excentricität der Ellipse abhängige. sondern allgemeine, Bestimmung der Anziehung, welche ein elliptischer Ring von unendlich kleiner und nach obigem
\SourceLeaf{367}
Gesetze unveränderlicher Dicke gegen einen jeden Punkt im Raume ausübt, ist
daher für die physische Astronomie von hohem Interesse. Inzwischen ist sie nicht
weniger auch in rein mathematischer Hinsicht merkwürdig, wegen der mancherlei Kunstgriffe, welche ihre vollständige Auflösung erfordert.

Von der Auflösung selbst ist es nicht wohl thunlich, hier einen Auszug zu
geben. Der Verf. hat eine rein analytische Behandlung gewählt; Kenner, welche
ihr mit Aufmerksamkeit folgen, werden leicht die geometrischen Correlate der
einzelnen in der Untersuchung vorkommenden Grössen, und die Umschmelzbarkeit in eine geometrische Form wahrnehmen. Hier mag es genügen, nur das
Endresultat anzuführen. Drei unbekannte Grössen $G,\ G',\ G''$ werden durch
die Wurzeln einer cubischen Gleichung bestimmt, aus deren Beschaffenheit sich
beweisen lässt, dass sie alle Mal drei reelle Wurzeln habe. Die nach einer beliebigen Richtung zerlegte Anziehung des elliptischen Ringes wird sodann durch
einen Ausdruck von der Form $pP+qQ$ dargestellt, wo $p$ und $q$ algebraisch
von $G,\ G',\ G''$ abhängen, $P$ und $Q$ hingegen die bestimmten Werthe der Integrale
\[
\int\frac{\cos T^{2}\,.\,\mathrm{d}T}
{2\pi(mm\cos T^{2}+nn\sin T^{2})^{\frac{3}{2}}}
\]
\[
\int\frac{\sin T^{2}\,.\,\mathrm{d}T}
{2\pi(mm\cos T^{2}+nn\sin T^{2})^{\frac{3}{2}}}
\]
bedeuten, wenn die Integrationen von $T=0$ bis $T=360^{0}=2\pi$ ausgedehnt
werden, und wo
\[
m=\surd(G+G'),\quad n=\surd(G+G'')
\]
\par Da diese Integrale ($m=n$ ausgenommen) transscendenter Natur sind. und bekannter Massen mit andern in der Perturbationsrechnung vorkommenden vielbehandelten Transscendenten zusammenhängen, so konnte die Auflösung, nachdem
sie bis auf diesen Punkt geführt war, als vollendet angesehen werden. Der Verfasser hat indessen diese erste sich ihm darbietende Gelegenheit benutzt, um die
ersten Linien eines \textit{neuen Algorithmus} zu geben, dessen er sich schon seit einer
langen Reihe von Jahren zur Bestimmung dieser Transscendenten bedient hat, und
worüber er in Zukunft eine ausgedehnte zu vielen merkwürdigen Resultaten führende Untersuchung bekannt machen wird. Hier können nur die Hauptsätze,
mit Uebergehung der Beweise angeführt werden. Wenn man aus zwei gegebe\SourceLeaf{368}%
nen positiven Grössen $m$ und $n$, andere $m',\ m'',\ m'''$ u. s. w., $n',\ n'',\ n'''$ u. s. w. nach
folgenden Gesetzen ableitet:
\[
\begin{array}{ll}
m'=\frac{1}{2}(m+n),&n'=\surd mn\\[0.4em]
m''=\frac{1}{2}(m'+n'),&n''=\surd m'n'\\[0.4em]
m'''=\frac{1}{2}(m''+n''),&n'''=\surd m''n''\ \mbox{u. s. w.}
\end{array}
\]
d. i. wenn $m',\ n'$ resp. das arithmetische und geometrische Mittel zwischen $m$
und $n$ ist; eben so $m'',\ n''$ das arithmetische und geometrische Mittel zwischen
$m'$ und $n'$ u. s. f.: so nähern sich die Glieder sowohl der Reihe $m,\ m',\ m'',\ m'''$ u. s. w.,
als die der Reihe $n,\ n',\ n'',\ n'''$ u. s. w. äusserst schnell einer gemeinschaftlichen
Grenze $=\mu$, welche der Verfasser das arithmetisch-geometrische Mittel von
$m$ und $n$ nennt. Offenbar ist $\mu$ zugleich das arithmetisch-geometrische Mittel
von $m'$ und $n'$, oder überhaupt von je zweien zusammengehörigen Gliedern der
beiden Reihen. Der Verfasser beweist nun, dass $\frac{1}{\mu}$ der Werth des Integrals
\[
\int\frac{\mathrm{d}T}{2\pi\surd(mm\cos T^{2}+nn\sin T^{2})}
\]
ist, wenn die Integration von $T=0$ bis $T=360^{0}$ ausgedehnt wird. Man
wird leicht sehen, dass diess auch auf folgende Art hätte ausgesprochen werden
können: Wenn die Entwicklung der Function
\[
\frac{1}{\surd(\alpha+\beta\cos\psi)}
\]
die Reihe
\[
A+B\cos\psi+C\cos2\psi+D\cos3\psi+\mathrm{etc.}
\]
gibt, so ist allezeit $\frac{1}{A}$ das arithmetisch-geometrische Mittel der beiden Grössen
$\surd(\alpha+\beta)$ und $\surd(\alpha-\beta)$.

Ein zweites eben so wichtiges Theorem ist, dass wenn man die Summe der
unendlichen jederzeit sehr schnell convergirenden Reihe
\[
2(m'm'-n'n')+4(m''m''-n''n'')+8(m'''m'''-n'''n''')+\mbox{u. s. w.}
\]
wo die Zahlencoëfficienten eine geometrische Progression bilden, $=(mm-nn)\nu$
setzt, der Werth des Integrals
\[
\int\frac{\cos2T\,.\,\mathrm{d}T}
{2\pi\surd(mm\cos T^{2}+nn\sin T^{2})}
\]
\SourceLeaf{369}
von $T=0$ bis $T=360^{0}$ erstreckt, $=-\frac{\nu}{\mu}$ wird. Offenbar ist denn hierdurch auch der zweite Coëfficient obiger Reihe bekannt, nemlich $B=-\frac{\nu}{2\mu}$
wenn man $m=\surd(\alpha+\beta)$, $n=\surd(\alpha-\beta)$ gesetzt hat. Alle folgenden Coëfficienten $C$, $D$, u.\,s.\,w. aber werden bekanntlich durch die beiden ersten $A$ und $B$
algebraisch und einfach bestimmt. Für die numerische Berechnung der Grössen
$m'm'-n'n'$, $m''m''-n''n''$ u.\,s.\,w. wird in der Abhandlung selbst noch ein besonderes
sehr bequemes Verfahren gelehrt.
\par
Die Anwendung auf die Transscendenten der gegenwärtigen Untersuchung
gibt endlich noch die einfachen Ausdrücke
\[
P=\frac{1+\nu}{2mm\mu},\qquad Q=\frac{1-\nu}{2nn\mu}
\]
\par Aufmerksamen Lesern wird es nicht entgehen, wie viele interessante Aufgaben, die mit den hier betrachteten Transscendenten zusammenhangen, durch
den erklärten Algorithmus mit grösster Leichtigkeit aufgelöst werden. Als ein
Beispiel führen wir hier die Rectification der Ellipse an. Setzt man ihre halbe
grosse Axe $=m$, die halbe kleine Axe $=n$, so wird die Peripherie
\[
\begin{aligned}&=\frac{2\pi}{\mu}\bigl\{m'm'-2(m''m''-n''n'')\bigr.\\[5pt]&\qquad\bigl.-4(m'''m'''-n'''n''')-8(m''''m''''-n''''n'''')-\mathrm{u.\,s.\,w.}\bigr\}\end{aligned}
\]
Ein anderes Beispiel gibt die Dauer der Pendelschwingungen bei endlichen Bogen, welche sich zu der Dauer der unendlich kleinen Schwingungen verhält, wie
die Einheit zu dem arithmetisch-geometrischen Mittel zwischen $1$ und dem Cosinus von einem Viertel des ganzen Schwingungsbogens.
\par
Schliesslich mus noch bemerkt werden, dass der Verf. diese Resultate, so
wie er sie schon vor vielen Jahren unabhängig von ähnlichen Untersuchungen
\textsc{Lagrange}'s und \textsc{Legendre}'s gefunden hat, in ihrer ursprünglichen Form darstellen
zu müssen geglaubt hat, obgleich sie zum Theil aus den Entdeckungen dieser Geometer leicht hätten abgeleitet werden können, theils weil jene Form ihm wesentliche Vorzüge zu haben schien, theils weil sie gerade so den Anfang einer viel
ausgedehntern Theorie ausmachen, wo seine Arbeit eine ganz verschiedene Richtung von der der genannten Geometer genommen hat.
\begin{center}\rule{0.25\linewidth}{0.4pt}\end{center}
\clearpage\fancyhead[L]{\small ARITHMETISCH GEOMETRISCHES MITTEL}
\pdfbookmark[0]{ARITHMETISCH GEOMETRISCHES MITTEL}{work-native370}
\SourceLeaf{370}
\selectlanguage{english}
\begin{center}
\textbf{NACHLASS.}\par
\textbf{[ARITHMETISCH GEOMETRISCHES MITTEL.]}\par
\textbf{PARS I.}\par
\textbf{DE ORIGINE PROPRIETATIBUSQUE}\\\textbf{GENERALIBUS NUMERORUM}\par
\textbf{MEDIORUM ARITHM. GEOMETRICORUM.}\par
\rule{0.09\linewidth}{0.4pt}\par
\end{center}
\AGMArticle{1}
Sint $\left\{\begin{array}{l}a,\allowbreak \ a',\allowbreak \ a'',\allowbreak \ a'''\ .\,\allowbreak .\,\allowbreak .\,\allowbreak .\\ b,\allowbreak \ b',\allowbreak \ b'',\allowbreak \ b'''\ .\,\allowbreak .\,\allowbreak .\,\allowbreak .\end{array}\right\}$ duae progressiones quantitatum ea lege formatae, ut
quilibet ipsarum termini correspondentes sint \textit{media} inter terminos antecedentes,
et quidem termini progressionis superioris media arithmetica, progressionis inferioris geometrica, puta
\[
\begin{aligned}&a'=\tfrac12(a+b),\quad b'=\surd ab,\\[6pt]&a''=\tfrac12(a'+b'),\quad b''=\surd a'b',\\[6pt]&a'''=\tfrac12(a''+b''),\quad b'''=\surd a''b''\ \mathrm{etc.}\end{aligned}
\]
Supponemus autem, ipsos $a$, $b$ esse reales positivos, et pro radicibus quadraticis
ubique accipi valores positivos; quo pacto progressiones quousque libuerit produci
poterunt, omnes ipsarum termini erunt plene determinati valoresque positivos reales nanciscentur. Porro prima hic se fronte offerunt observationes sequentes:
\par
I. Si $a=b$, omnes utriusque seriei termini erunt $=a=b$.
\par
II. Si vero $a$, $b$ sunt inaequales, erit $(a'-b')(a'+b')=\tfrac14(a-b)^2$, unde
concluditur $b'<a'$, et perinde erit $b''<a''$, $b'''<a'''$ etc., i.\,e. quivis terminus seriei inferioris minor erit quam correspondens superioris. Quocirca in hoc casu
supponemus, esse etiam $b<a$.
\par
III. Eadem suppositione erit $a'<a$, $b'>b$; $a''<a'$, $b''>b'$ etc.; progressio
itaque superior continuo decrescit, inferior continuo crescit; hinc manifestum est,
utramque habere limitem; hi limites commode exprimuntur per $a^{\infty}$, $b^{\infty}$.
\par
IV. Denique ex $\frac{a'-b'}{a-b}=\frac{(a-b)}{4(a'+b')}=\frac{a-b}{2(a+b)+4b'}$ sequitur $a'-b'<\tfrac12(a-b)$,
eodemque modo erit $a''-b''<\tfrac12(a'-b')$ etc. Hinc concluditur, $a-b$, $a'-b'$,
\SourceLeaf{371}
$a''-b''$, $a'''-b'''$ etc. constituere progressionem continuo decrescentem atque ipsius
limitem esse $=0$. Hinc $a^{\infty}=b^{\infty}$, i.\,e. progressio superior et inferior eundem
limitem habebunt, quo illa semper manet maior, haec minor.
\par
Hunc limitem vocamus \textit{numerum medium arithmetico-geometricum inter $a$ et $b$},
et per $\mathrm{M}(a,b)$ designamus.
\AGMArticle{2}
Radices aequationis $xx-2ax+bb=0$ erunt reales positivi, siquidem
$a\geqq b$; medium arithmeticum inter has radices erit $a$, geometricum $b$; designata
itaque una radice (et quidem maiore si sunt inaequales) per ${}^{\prime}\!a$, altera per ${}^{\prime}\!b$, poterit ${}^{\prime}\!a$ spectari tamquam terminus progressionis superioris terminum $a$ praecedens, eodemque modo ${}^{\prime}\!b$ tamquam terminus progressionis inferioris ante $b$. Similiter designando
\begingroup\setlength{\arraycolsep}{2pt}
\[
\begin{array}{@{}lllll@{}}
\mathrm{aequationis}&xx-2\,{}^{\prime}\!ax+{}^{\prime}\!b\,{}^{\prime}\!b=0&\mathrm{radicem\ maiorem\ per}&{}^{\prime\prime}\!a&\mathrm{minorem\ per}\ {}^{\prime\prime}\!b\\[6pt]
&xx-2\,{}^{\prime\prime}\!ax+{}^{\prime\prime}\!b\,{}^{\prime\prime}\!b=0&&{}^{\prime\prime\prime}\!a&\hfill{}^{\prime\prime\prime}\!b\\[6pt]
&xx-2\,{}^{\prime\prime\prime}\!ax+{}^{\prime\prime\prime}\!b\,{}^{\prime\prime\prime}\!b=0&&{}^{\prime\prime\prime\prime}\!a&\hfill{}^{\prime\prime\prime\prime}\!b
\end{array}
\]\endgroup
poterunt ${}^{\prime}\!a$, ${}^{\prime\prime}\!a$, ${}^{\prime\prime\prime}\!a$ etc. spectari tamquam continuatio progressionis superioris versus laevam, atque ${}^{\prime}\!b$, ${}^{\prime\prime}\!b$, ${}^{\prime\prime\prime}\!b$ etc. tamquam continuatio progressionis inferioris, ita
ut iam habeantur duae progressiones utrimque in infinitum continuabiles
\[
\begin{array}{@{}ll@{}}
.\,.\,.\,.\ {}^{\prime\prime\prime\prime}\!a,\;{}^{\prime\prime\prime}\!a,\;{}^{\prime\prime}\!a,\;{}^{\prime}\!a,\; a,\; a',\; a'',\; a''',\; a''''\ .\,.\,.\,.\,.&\mathrm{(I)}\\[6pt]
.\,.\,.\,.\ {}^{\prime\prime\prime\prime}\!b,\;{}^{\prime\prime\prime}\!b,\;{}^{\prime\prime}\!b,\;{}^{\prime}\!b,\; b,\; b',\; b'',\; b''',\; b''''\ .\,.\,.\,.\,.&\mathrm{(II)}
\end{array}
\]
\par Quivis itaque terminus progressionis (I) erit maior quam correspondens seriei (II); series illa a laeva ad dextram continuo decrescit, a dextra ad laevam
crescit: haec a laeva ad dextram continuo crescit sensuque contrario decrescit.
Versus dextram utraque series eundem limitem habet; versus laevam autem (I)
super omnes limites crescit, (II) habet limitem $0$ (nisi omnes utriusque progressionis termini sunt aequales). Nam ${}^{\prime}\!a=a+\surd(aa-bb)$; ${}^{\prime}\!b=a-\surd(aa-bb)$;
hinc ${}^{\prime}\!a\,{}^{\prime}\!a-{}^{\prime}\!b\,{}^{\prime}\!b=4a\surd(aa-bb)>4(aa-bb)$, similiterque ${}^{\prime\prime}\!a\,\allowbreak {}^{\prime\prime}\!a-{}^{\prime\prime}\!b\,\allowbreak {}^{\prime\prime}\!b>4({}^{\prime}\!a\,\allowbreak {}^{\prime}\!a-{}^{\prime}\!b\,\allowbreak {}^{\prime}\!b)$
etc., unde patet, seriem $aa-bb$, ${}^{\prime}\!a\,{}^{\prime}\!a-{}^{\prime}\!b\,{}^{\prime}\!b$, ${}^{\prime\prime}\!a\,{}^{\prime\prime}\!a-{}^{\prime\prime}\!b\,{}^{\prime\prime}\!b$ etc. et proin etiam hanc
$a$, ${}^{\prime}\!a$, ${}^{\prime\prime}\!a$ etc. quemvis limitem superare posse; $\frac{{}^{\prime}\!b}{b}=\frac{b}{{}^{\prime}\!a}<\frac{b}{a}$ similiterque $\frac{{}^{\prime\prime}\!b}{{}^{\prime}\!b}<\frac{{}^{\prime}\!b}{{}^{\prime}\!a}$ etc.
i.\,e. $\frac{{}^{n+1}\!b}{{}^{n}\!b}$ infra quemvis limitem deprimi potest augendo ipsum $n$, adeoque limes
seriei $b$, ${}^{\prime}\!b$, ${}^{\prime\prime}\!b$ etc. $=0$.
